% day2_seg3_kalstats/cov_mat.tex
\begin{frame}
  \onslide<1->{
	\frametitle{Covariance Matrices} 
	\framesubtitle{$Q,R,P$}
}
\begin{itemize}
  \onslide<2->{
 \item There are several covariance matrices associated with the Kalman filter: 
 $ \procnoisemat{k},\measnoisemat{k}, P(k)$
}
  \onslide<3->{
 \item In general, a covariance matrix is associated with a random vector $X$ as in the following definition:
\[C_X = \expect{XX\trans}\]
}
\end{itemize}
\end{frame}
\begin{frame}
  \onslide<1->{
	\frametitle{Covariance Matrices} 
	\framesubtitle{$Q,R,P$}
}
\begin{itemize}
  \onslide<2->{
 \item This is precisely how $\procnoisemat{k}$ and $\measnoisemat{k}$ are defined:
\[\procnoisemat{k}=\expect{w(k)w\trans(k)}\] and
\[\measnoisemat{k}=\expect{v(k)v\trans(k)}\]
}
  \onslide<3->{
\item $\covmat{k}{+}$, on the other hand is called the error covariance, because it is associated with $\esterr$
\[\covmatdef\]
where $\esterrdef$.
}
\end{itemize}
\end{frame}
\begin{frame}
  \onslide<1->{
	\frametitle{Covariance Matrices} 
	\framesubtitle{Positive Definite Symmetric Matrices}
}
\begin{itemize}
  \onslide<2->{
\item Covariance Matrices are \textbf{positive definite}: The quadratic form $x\trans C x>0$ for all nonzero vectors $x$  
}
  \onslide<3->{
\item We write $C>0$ for positive definite matrices  
}
  \onslide<4->{
\item $C>0$ iff all eigenvalues of $C$ are positive 
}
  \onslide<5->{
\item All diagonal elements of $C$ are positive if $C>0$  
}
  \onslide<6->{
\item If $C>0$ and $S$ non-singular (invertible), then $SCS\trans>0$ and $S\trans CS>0$
}
  \onslide<7->{
\item If $C>0$ then $C$ is non-singular and $C^{-1}>0$
}
\end{itemize}
\end{frame}

\begin{frame}
  \onslide<1->{
	\frametitle{Covariance Matrices} 
	\framesubtitle{Lower Triangular Cholesky Decomposition}
}
\begin{itemize}
  \onslide<2->{
\item  If $C>0$ then a lower triangular matrix $L$ can be found such that $C=LL\trans$ 
}
  \onslide<3->{
\item $L$ is called a \textbf{matrix square root} of $C$  
}
  \onslide<4->{
\item  Square root matrices are not unique;  $TL$ is also a square root of $C$ if $L$ is such and if $T$ is orthogonal $TT\trans=I$ 
}
  \onslide<5->{
\item We'll learn more about square root matrices when we study square root filtering  
}
\end{itemize}
\end{frame}
